\documentclass{article}

\usepackage{PRIMEarxiv}

\usepackage[utf8]{inputenc}
\usepackage[T1]{fontenc}
\usepackage{hyperref}
\usepackage{url}
\usepackage{booktabs}
\usepackage{amsfonts}
\usepackage{amsmath}
\usepackage{amssymb}
\usepackage{nicefrac}
\usepackage{microtype}
\usepackage{fancyhdr}
\usepackage{graphicx}
\usepackage{multirow}
\usepackage{array}
\usepackage{xcolor}
\usepackage{enumitem}
\usepackage{caption}
\usepackage{subcaption}
\usepackage{natbib}        % for \citep / \citet
\graphicspath{{media/}}

% ── Header ───────────────────────────────────────────────────────────────────
\pagestyle{fancy}
\thispagestyle{empty}
\rhead{\textit{}}
\fancyhead[LO]{SpineSurg-CT: L6 Segmentation and LSTV Detection on CT}

% ── Placeholder command (orange bold — remove when results available) ─────────
\definecolor{placeholder}{RGB}{180,100,0}
\newcommand{\ph}[1]{\textcolor{placeholder}{\textbf{[#1]}}}

% ─────────────────────────────────────────────────────────────────────────────
\title{SpineSurg-CT: A Deep Learning Model for Automated Segmentation
and Detection of Lumbosacral Transitional Vertebrae on CT Imaging%
\thanks{\textit{\underline{Citation}}: Schwing G.
SpineSurg-CT: A Deep Learning Model for Automated Segmentation and
Detection of Lumbosacral Transitional Vertebrae on CT Imaging.
\textit{arXiv} preprint arXiv:XXXX.XXXXX, 2026.}}

\author{
  Gregory Schwing, MD, PhD \\
  Department of Surgery (Preliminary) \\
  Detroit Medical Center / Wayne State University \\
  Detroit, MI, USA \\
  \texttt{gregory.schwing@wayne.edu}
  \And
  Patrick J. Schwing, BS, MS \\
  LSU Health New Orleans School of Medicine \\
  New Orleans, LA 70112, USA \\
  \And
  Hazem Eltahawy, MD, PhD \\
  Department of Neurosurgery \\
  Detroit Medical Center \\
  Detroit, MI, USA \\
}

% ─────────────────────────────────────────────────────────────────────────────
\begin{document}
\maketitle

% ── Abstract ─────────────────────────────────────────────────────────────────
\begin{abstract}
\noindent\textbf{Purpose:} To develop and evaluate SpineSurg-CT, the first
CT segmentation model with lumbosacral transitional vertebra (LSTV) awareness
built into its architecture --- providing voxel-level spatial segmentation of
the supernumerary sixth lumbar vertebra (L6) and a simultaneous probabilistic
classification of LSTV subtype (Normal / Lumbarization / Sacralization) in a
single end-to-end inference pass from raw CT, without requiring upstream
vertebra localization or counting.

\noindent\textbf{Materials and Methods:} In this retrospective study, 784 CT
examinations from the publicly available CTSpine1K COLONOG dataset were used
for training ($n{=}552$), validation ($n{=}116$), and testing ($n{=}116$) with
stratified sampling ensuring proportional LSTV representation across all
partitions. The cohort included 16 lumbarization cases (L6 present), 9
sacralization cases (L5 fused to sacrum), and 759 normal spines. A 3D ResUNet
with cross-scale feature fusion modules (CTSpineSegNet architecture) was
extended with: (1)~a 10-class segmentation output with an explicit L6 class
under direct voxel-level supervision; (2)~a secondary LSTV classification head
at the bottleneck producing calibrated Normal/Lumbarization/Sacralization
probabilities; and (3)~a multi-strategy class imbalance mitigation stack
including novel registration-based sacralization synthesis using ANTs SyN
deformation fields computed and applied entirely in preprocessed
$256^3$ unit coordinate space. Performance was evaluated using Dice similarity
coefficient (DSC), 95th-percentile Hausdorff distance (HD95), and mean
intersection over union (mIOU), benchmarked against TotalSegmentator.

\noindent\textbf{Results:} \ph{SpineSurg-CT achieved mean L6 DSC of X.XX
(95\% CI: X.XX--X.XX) on lumbarization test cases ($n{=}3$), compared with
0.00 for TotalSegmentator. LSTV subtype classification yielded
Normal/Lumbarization/Sacralization probabilities with overall accuracy X.XX.
Mean L1--L5 DSC was X.XX. All comparisons against TotalSegmentator on LSTV
cases were statistically significant (Wilcoxon signed-rank, $p < 0.001$).}

\noindent\textbf{Conclusion:} SpineSurg-CT is the first CT segmentation model
with LSTV awareness built into its architecture, simultaneously delivering a
voxel-level 3D L6 mask and calibrated LSTV subtype probabilities without
upstream localization. This represents a clinically meaningful advance for
preventing wrong-level spine surgery: existing tools either segment without
LSTV awareness (TotalSegmentator), label without segmenting (VERIDAH), or
classify without CT spatial output (Kwak et al.).
\end{abstract}

\keywords{lumbosacral transitional vertebra \and spine segmentation \and
computed tomography \and deep learning \and wrong-level surgery \and L6 \and
lumbarization \and sacralization}

% ── Front-matter boxes ───────────────────────────────────────────────────────
\noindent\textbf{Summary Statement:} SpineSurg-CT is the first CT segmentation
model with lumbosacral transitional vertebra (LSTV) awareness built into its
architecture, simultaneously delivering explicit voxel-level L6 segmentation and
calibrated Normal/Lumbarization/Sacralization probability output in a single
end-to-end pass from raw CT --- without requiring upstream vertebra localization
or vertebra counting.

\vspace{0.5em}
\noindent\textbf{Key Results:}
\begin{enumerate}[leftmargin=*,label=\arabic*.,noitemsep]
  \item \ph{SpineSurg-CT achieved mean L6 DSC of X.XX on lumbarization test
    cases, compared with 0.00 for TotalSegmentator, which lacks an L6 class
    by architectural design.}
  \item \ph{The LSTV classification head produced calibrated
    Normal/Lumbarization/Sacralization probabilities with F1 of
    X.XX/X.XX/X.XX --- the first probabilistic multi-class LSTV classifier
    embedded in a CT segmentation architecture, operating without upstream
    localization or vertebra counting.}
  \item \ph{Mean L1--L5 DSC was X.XX, comparable to TotalSegmentator (X.XX),
    with superior lumbosacral junction performance in LSTV cases.}
\end{enumerate}

\vspace{0.5em}
\noindent\textbf{Abbreviations:} CT, computed tomography; DSC, Dice similarity
coefficient; HD95, 95th-percentile Hausdorff distance; LSTV, lumbosacral
transitional vertebra; mIOU, mean intersection over union; L6, supernumerary
sixth lumbar vertebra; CFFM, cross-scale feature fusion module.

% =============================================================================
\section{Introduction}
% =============================================================================

Wrong-level spine surgery is a sentinel event in neurosurgery, classified as a
never-event under the Universal Protocol. Large surgical series consistently
report incidence below 1\%, but given the volume of lumbar procedures performed
annually this represents a substantial preventable burden of harm with severe
clinical and medicolegal consequences~\citep{albano2024}. Among the anatomical
risk factors implicated in vertebral level miscounting, lumbosacral transitional
vertebrae (LSTV) are consistently identified as a primary
precipitant~\citep{albano2024,karsy2024}. In a prospective series of 550
consecutive lumbar surgery patients, Ishida et al.\ documented LSTV in 12.9\%
of cases, with vertebral miscounting at initial outpatient evaluation in 54\% of
LSTV cases and surgical level revision after intraoperative counting in
15\%~\citep{ishida2017}.

LSTV encompasses two anatomical subtypes with opposing counting implications.
In \textit{lumbarization}, the uppermost sacral segment fails to fuse with the
remainder of the sacrum and develops as a supernumerary sixth lumbar vertebra
(L6), increasing the mobile segment count by one. In \textit{sacralization},
the fifth lumbar vertebra fuses partially or completely with the sacrum, reducing
the mobile segment count by one. LSTV prevalence in the surgical population is
4--18\% depending on imaging modality and classification
criteria~\citep{apazidis2011,nardo2012}.

Computed tomography is the reference standard for LSTV detection and
classification. Hanhivaara et al.\ demonstrated that CT achieves 76\%
sensitivity and 93\% specificity for LSTV, compared with 54\% sensitivity for
standard lumbar MRI without coronal sequences~\citep{hanhivaara2024}. Automated
CT segmentation could serve as a systematic screening step when CT is available,
or as a confirmation tool when LSTV is suspected.

No existing automated CT segmentation system can detect LSTV or label L6.
TotalSegmentator --- the most widely deployed general-purpose CT segmentation
tool, trained on 1,204 CT examinations segmenting 104 anatomical structures ---
assumes a fixed five-lumbar anatomy and has no L6 class by
design~\citep{wasserthal2023}. VERIDAH~\citep{moller2026} is a sophisticated vertebra labeling
algorithm achieving 97.22\% lumbar enumeration anomaly recall on CT. It uses a
3D DenseNet169 with four learned classification heads (vertebra label, anatomical
region, transition type, visibility) and a constrained sequence optimizer ---
it does learn morphological features from CT. However, L6 is represented
\textit{implicitly}: during training, L6 labels are replaced with two consecutive
L5 labels, and the sequence optimizer detects lumbarization by finding that two
consecutive L5-class predictions minimise the global path cost. There is no
explicit L6 class, no direct voxel-level L6 supervision, and no 3D spatial mask.
VERIDAH outputs vertebra centroid labels and requires pre-computed vertebra
localizations from an upstream segmentation model (SPINEPS), precluding
end-to-end inference from raw CT

Critically, Schinz et al.\ recently demonstrated at the \textit{thoracolumbar}
junction that morphometric vertebral shape-based classification outperforms
counting-based approaches on all metrics --- achieving 100\% agreement with
nerve morphology (the functional ground truth) versus 92.6\% for the
counting-based standard --- and is robust to incomplete spine imaging where
counting fails entirely~\citep{schinz2025}. This validates the principle that
morphological feature learning is superior to counting for vertebral transition
detection. However, Schinz et al.\ provide a manual classification scheme
applied by expert radiologists; no automated neural network model for this task
exists, and the lumbosacral junction --- clinically more consequential for
wrong-level surgery --- was explicitly not assessed.

A ResNet-50 classifier operating on plain radiographs achieves 76.4\% accuracy
for binary LSTV detection but provides no segmentation, no subtype distinction,
and requires a separate imaging modality~\citep{kwak2025}.

We therefore developed SpineSurg-CT, extending
CTSpineSegNet~\citep{li2025spinenet} with three novel components for LSTV
detection: (1)~a 10-class segmentation output including an explicit L6 class,
(2)~a secondary multi-task LSTV classification head predicting
Normal/Lumbarization/Sacralization from learned morphological features rather
than counting, and (3)~a multi-strategy class imbalance mitigation stack
including registration-based sacralization synthesis. To our knowledge, SpineSurg-CT is the first CT segmentation model
with an explicit L6 class under direct voxel-level supervision, producing a 3D
spatial segmentation mask without upstream localization steps. Whereas VERIDAH
detects lumbarization through implicit sequence optimization over relabelled L5
centroids, SpineSurg-CT learns L6 morphology directly from annotated voxels,
extending the morphometric classification principle validated by Schinz et al.\
at the thoracolumbar junction~\citep{schinz2025} to a fully automated spatial
segmentation framework at the lumbosacral junction.

% =============================================================================
\section{Materials and Methods}
% =============================================================================

\subsection{Study Design}

This retrospective study used publicly available datasets. Institutional review
board approval and informed consent were not required. The study was conducted
and reported in accordance with the Checklist for Artificial Intelligence in
Medical Imaging (CLAIM) 2024 guidelines~\citep{claim2024}.

\subsection{Datasets}

\subsubsection{CTSpine1K --- Primary Training Dataset}

CTSpine1K~\citep{li2021ctspine1k} is a publicly available vertebral
segmentation benchmark comprising 1,005 CT examinations across four clinical
subsets. The COLONOG subset (784 cases, colonoscopy screening CTs) was selected
as the appropriate training cohort for lumbosacral segmentation: all cases
include complete L1--L5 lumbar coverage in the supine position. The COVID-19
and HNSCC-3DCT subsets were excluded due to truncated field of view, confirmed
by absence of complete lumbar labels in a pre-training audit.

Within the 784 COLONOG cases, LSTV prevalence was determined by systematic
label scanning: 16 cases contained raw label 25 (L6, lumbarization) and 9
cases had raw label 24 (L5) absent with remaining lumbar labels present
(sacralization), for 25 total LSTV cases (3.2\% prevalence). A unified 10-class
label mapping was applied: background (0), L1--L5 (1--5), L6 (6), sacrum (7),
left hip (8), right hip (9).

\subsubsection{CTPelvic1K --- Supplementary Pelvic Supervision}

CTPelvic1K~\citep{liu2021ctpelvic1k} Dataset~1 contains pelvic masks for 35
COLONOG cases. Datasets 2--7 originate from independent clinical studies with
no patient overlap with CTSpine1K COLONOG, confirmed by DICOM UID analysis.
Patient-number-based matching of Dataset~1 to CTSpine1K identified 33/35
matched pairs (94\% match rate). No LSTV cases overlap with the 33 matched
pelvic cases (zero intersection confirmed empirically). The 33 matched cases
provide sacrum (class 7), left hip (class 8), and right hip (class 9) ground
truth via online nearest-neighbour resampling into spine CT space during
training. Critically, these 33 cases are the only samples in the dataset with
complete lumbosacral label coverage including the pelvic girdle; the remaining
751 cases carry lumbar vertebra labels only.

\subsection{Dataset Split}

Dataset partitioning used stratified sampling across three priority-ordered
strata: (1)~LSTV cases ($n{=}25$); (2)~normal spine cases with matched pelvic
mask ($n{=}33$); (3)~normal spine-only cases ($n{=}726$). Each stratum was
independently shuffled (seed = 42) and split 70/15/15, yielding:
train $n{=}552$, validation $n{=}116$, test $n{=}116$. This guarantees a
minimum of 1 LSTV case in both validation and test sets, and ensures that
pelvic mask cases appear in all three partitions.
Table~\ref{tab:dataset} summarises dataset characteristics.

\begin{table}[ht]
\centering
\caption{Dataset Characteristics by Stratum and Partition}
\label{tab:dataset}
\begin{tabular}{lrrrrr}
\toprule
Stratum & Total & Train & Val & Test & LSTV \\
\midrule
LSTV cases                       & 25  & 17  & 4   & 4   & 25 (100\%) \\
\quad Lumbarization (L6 present) & 16  & 11  & 2   & 3   & ---        \\
\quad Sacralization (L5 fused)   &  9  &  7  & 1   & 1   & ---        \\
Spine + pelvic mask              & 33  & 23  & 5   & 5   & 0          \\
Spine only                       & 726 & 512 & 107 & 107 & 0          \\
\midrule
\textbf{Total}                   & \textbf{784} & \textbf{552} &
  \textbf{116} & \textbf{116} & 25 (3.2\%) \\
\bottomrule
\end{tabular}\\[2pt]
\footnotesize{LSTV, lumbosacral transitional vertebra.
Spine + pelvic mask cases have confirmed zero overlap with LSTV cases.
Stratified 70/15/15 split applied independently within each stratum.}
\end{table}

\subsection{Image Preprocessing}

All images were reoriented to PIR convention (posterior-inferior-right axial
orientation, per \citealt{li2025spinenet}), resampled to 1\,mm isotropic
spacing (bilinear for images, nearest-neighbour for labels), HU-clipped to
$[-1000, 1000]$ and linearly normalised to $[0,1]$, foreground-cropped based
on image intensity content (20-voxel margin), zero-padded if necessary, and
centre-cropped to $256^3$ voxels.

Image-intensity-based foreground cropping --- rather than label-based cropping
--- was used deliberately to preserve the full lumbosacral and pelvic anatomy
within the $256^3$ field of view. Label-based cropping computes the bounding
box from the spine labels only, which terminate at the inferior sacrum; the hip
bones (classes 8--9) extend 40--60\,mm below this boundary and would be
systematically excluded before pelvic mask fusion, rendering pelvic ground
truth absent from training despite being available. Image-based cropping
captures the full bony extent of the body in the field of view, ensuring
sacrum and hip labels are physically present in the training volume after
\texttt{FuseSpinePelvicLabels}. For cases without a matched pelvic mask, the
additional inferior region contains only background-class air voxels (intensity
$\approx 0$ after normalisation), which are indistinguishable from zero-padding
and impose no learning cost. This single preprocessing function is shared
identically between training, registration-cache computation, and 3D
augmentation visualisation, guaranteeing that what is visualised in the
quality-control output is exactly what the network receives during training.
The resulting volumes occupy a shared unit coordinate space (origin at zero,
1\,mm isotropic spacing, identity direction cosines) exploited by the
registration-based augmentation pipeline described in
Section~\ref{sec:class_imbalance}.

\subsection{Model Architecture}

SpineSurg-CT extends CTSpineSegNet~\citep{li2025spinenet}, a 3D residual U-Net
with cross-scale feature fusion modules (CFFM) at encoder skip connections
between layers $1{\to}2$ and $2{\to}3$, and a linear self-attention bottleneck
reducing complexity from $O(n^2)$ to $O(kn)$ with $k{=}64$. Three extensions
were added:

\begin{enumerate}[noitemsep]
  \item \textbf{10-class segmentation output:} L6 (class 6) added with a
    $20\times$ loss weight in the combined Dice-CE objective to compensate for
    extreme rarity (16/784 cases, 2.0\%).
  \item \textbf{LSTV classification head:} Secondary branch at the bottleneck:
    global average pooling $\to$ FC$(512{\to}128{\to}3)$ with ReLU and 0.3
    dropout, producing Normal/Lumbarization/Sacralization probabilities
    (loss weight $\alpha{=}0.5$).
  \item \textbf{Data augmentation:} Random flip ($p{=}0.5$, axis 0), rotation
    $\pm30^{\circ}$ ($p{=}0.5$, all axes), zoom $0.7$--$1.4\times$ ($p{=}0.5$),
    Gaussian noise $\sigma{\le}0.1$ ($p{=}0.3$), contrast jitter
    $\gamma{=}0.7$--$1.5$ ($p{=}0.3$), intensity scale/shift $\pm10\%$
    ($p{=}0.3$ each).
\end{enumerate}

Total trainable parameters: 15,877,165.

\subsection{Class Imbalance Mitigation}
\label{sec:class_imbalance}

Four strategies addressed LSTV class imbalance (3.2\% prevalence):

\paragraph{1. WeightedRandomSampler.}
LSTV cases oversampled at $10\times$, ensuring approximately one LSTV case per
batch at batch size 2.

\paragraph{2. Class-weighted loss.}
Per-class loss weights: $20\times$ (L6), $2\times$ (sacrum, hips), $1\times$
(L1--L5), $0.1\times$ (background). Sacralization samples receive an additional
$2\times$ sacrum weight boost applied per-sample in the cross-entropy term,
reflecting the anatomically enlarged sacrum produced by L5 incorporation.

\paragraph{3. CopyPaste lumbarization synthesis with targeted elastic deformation.}
The 11 lumbarization training cases were preprocessed offline into an
in-memory pool. During training, 30\% of normal spine samples received
synthetic lumbarization as follows. A random pool case was selected and its L6
label bounding box extracted with a 4-voxel margin. In lumbarization, L6 is the
supernumerary segment that interposes between L5 and the sacrum --- it is
\textit{inferior} to L5, occupying the L5--S1 disc space. Accordingly, the L6
patch was placed at the \textit{inferior} boundary of the target's L5 label along
the inferior-to-superior axis (axis 1 in PIR space): the patch was positioned
such that its superior face aligns with the inferior edge of target L5, reproducing
the anatomical relationship. Along the posterior-to-anterior axis (axis 0) and
right-to-left axis (axis 2), the patch was centred on the target spine column
centroid. Image intensities and label values were updated simultaneously within
the L6 mask region. Elastic pre-deformation of the source patch was deliberately
omitted: deforming the source before bounding box extraction scatters label voxels
outside the pre-computed box, producing floating label fragments disconnected
from the spine column.

Following the paste, targeted elastic deformation was applied exclusively to the
L5--L6--sacrum subvolume (bounding box of classes 5, 6, 7 with 24-voxel context
margin), leaving L1--L4 unperturbed. This is anatomically appropriate: in
lumbarization, the upper lumbar spine is completely normal and elastic deformation
of L1--L4 would introduce spurious morphological variation into the supervision
signal for these classes. The 24-voxel margin provides sufficient spatial context
to prevent displacement field discontinuities at the L4/L5 boundary without
deforming L4 itself. The identical displacement field was applied to both image
and label simultaneously using MONAI \texttt{Rand3DElasticD}: bilinear
interpolation for image intensities and nearest-neighbour interpolation for the
discrete label volume. The label requires no boundary blending --- discrete
segmentation masks have no intensity seam artefact --- while image intensities
at the subvolume boundary were smoothly composited with the original using a
Gaussian-feathered blend mask (linear fade over 24 voxels per face) to eliminate
visible intensity discontinuities at the elastic subvolume boundary.

\paragraph{4. Registration-based sacralization synthesis.}
This is the methodologically novel contribution of the present work. The core
challenge in sacralization augmentation is that simple label modification
(zeroing the L5 annotation) leaves CT image intensities unchanged, producing
supervision signal without a corresponding image-domain deformation. A
realistic synthetic sacralization case must simultaneously exhibit: collapsed
L5--S1 disc space, enlarged transverse processes, and superior elongation of
the sacrum.

\textit{Phase 1 --- offline registration (one-time cost, ${\approx}45$\,min).}
For each of the 7 sacralization cases in the training split, the CT volume was
preprocessed to $256^3$ 1\,mm isotropic and registered to a reference normal
spine using ANTs SyNRA deformable registration~\citep{avants2011}. Critically,
both registration and subsequent field application were performed entirely in
the shared preprocessed unit coordinate space (PIR, $256^3$, 1\,mm isotropic,
unit origin and direction), not in the original NIfTI physical space. This
eliminates the coordinate frame mismatch that arises when fields computed in
heterogeneous native NIfTI spaces are applied to standardised preprocessed
tensors. Reference normals were drawn exclusively from the 23 pelvic-masked
training cases, ensuring that the deformation target has complete sacrum and
hip labels (classes 7--9). The forward deformation field --- encoding the
geometric transformation from normal to sacralization morphology --- was cached
to disk alongside preprocessed volume arrays in NumPy format for instant
retrieval on subsequent runs.

\textit{Phase 2 --- online synthesis (per-batch, ${\approx}2$\,s).}
Unlike the CopyPaste augmentation, which applies targeted local elastic
deformation restricted to the L5--L6--sacrum subvolume, sacralization synthesis
applies the ANTs deformation field to the entire $256^3$ volume. This is
appropriate: the registration field was computed from a globally coherent
morphological difference between a real sacralization case and a normal reference,
encoding simultaneous changes across the entire lumbosacral anatomy --- caudal
migration of L5 toward the sacrum, collapse of the L5--S1 disc space, enlarging
transverse processes, and superior elongation of the sacrum. Truncating this field to a subvolume would arbitrarily corrupt a globally
consistent deformation and introduce artificial discontinuities absent in the
original anatomy.
During training, a matched normal case (one with a confirmed pelvic mask,
$n{=}23$ in training) was selected as the augmentation target. The Boolean
flag \texttt{has\_pelvic\_mask} was propagated through the preprocessing
pipeline to guarantee that field application is restricted to cases with
complete label coverage, regardless of MONAI tensor metadata state. A randomly
selected cached field was applied to both image and label tensors simultaneously
using ANTs linear interpolation (image) and nearest-neighbour (label). Post-warp
L5 voxels geometrically adjacent to the sacrum --- identified by binary dilation
of the sacrum mask with 8-voxel structuring element (8\,mm at 1\,mm isotropic)
--- were relabelled as sacrum (class 7), modelling L5 incorporation. If more
than 50\% of L5 was found adjacent to the sacrum, all remaining L5 voxels were
absorbed (complete sacralization). The synthetic label thus contains a fused
L5+sacrum structure with intact hip labels, reflecting realistic sacralization
anatomy. Because all 23 matched training cases are eligible and augmentation
probability is set to 1.0, approximately 23 anatomically complete synthetic
sacralization samples are generated per training epoch --- the hard ceiling
imposed by the number of pelvic-masked training cases in this dataset.
The ANTs deformation field was computed from spine anatomy only; its
magnitude decays smoothly to near-zero in the inferior pelvis region
where no corresponding spine structure drives the registration, so hip
label integrity in the synthetic output is preserved without explicit
masking or protection.

\subsection{Training Configuration}

Training used AdamW ($\text{lr}{=}10^{-3}$, weight decay $10^{-5}$), linear
warmup (5 epochs) then cosine annealing to $\eta_{\min}{=}10^{-6}$, maximum
300 epochs with early stopping (patience 30, monitor: \texttt{val/DSC},
$\Delta_{\min}{=}10^{-4}$), batch size 2, BF16 mixed precision, gradient
clipping 1.0. Hardware: NVIDIA H200 NVL (141\,GB VRAM, Hopper architecture).
\texttt{torch.set\_float32\_matmul\_precision(``high'')} was set to exploit
native BF16 tensor cores. Framework: PyTorch 2.4.1 (CUDA 12.4), PyTorch
Lightning 2.4.0, MONAI 1.4.0, antspyx (ANTs Python bindings).

\subsection{Evaluation}

Primary evaluation compared SpineSurg-CT against TotalSegmentator (v2.4.0,
CT model) on the 116-case held-out test set. Metrics: DSC, HD95 (mm), and
mIOU per class. LSTV classification: sensitivity, specificity, F1, and AUROC
per class. Statistical significance: Wilcoxon signed-rank test, $p{<}0.05$.
Confidence intervals: bootstrap resampling ($n{=}1000$). The pre-specified
\textbf{primary outcome} is L6 DSC on lumbarization test cases;
TotalSegmentator L6 DSC is 0.00 by design.

% =============================================================================
\section{Results}
% =============================================================================

\subsection{Segmentation Performance}

\ph{Table~\ref{tab:segmentation} reports per-class DSC, HD95, and mIOU for
SpineSurg-CT and TotalSegmentator on the 116-case test set. SpineSurg-CT
achieved mean L1--L4 DSC of X.XX\,$\pm$\,X.XX, comparable to TotalSegmentator
(X.XX\,$\pm$\,X.XX, $p{=}$X.XX).}

\begin{table}[ht]
\centering
\caption{Segmentation Performance on Held-Out Test Set ($n{=}116$)}
\label{tab:segmentation}
\begin{tabular}{lcccccc}
\toprule
 & \multicolumn{2}{c}{DSC} & \multicolumn{2}{c}{HD95 (mm)}
 & \multicolumn{2}{c}{mIOU} \\
\cmidrule(lr){2-3}\cmidrule(lr){4-5}\cmidrule(lr){6-7}
Class & SpineSurg & TotalSeg & SpineSurg & TotalSeg
      & SpineSurg & TotalSeg \\
\midrule
L1           & \ph{X.XX} & \ph{X.XX} & \ph{X.XX} & \ph{X.XX} & \ph{X.XX} & \ph{X.XX} \\
L2           & \ph{X.XX} & \ph{X.XX} & \ph{X.XX} & \ph{X.XX} & \ph{X.XX} & \ph{X.XX} \\
L3           & \ph{X.XX} & \ph{X.XX} & \ph{X.XX} & \ph{X.XX} & \ph{X.XX} & \ph{X.XX} \\
L4           & \ph{X.XX} & \ph{X.XX} & \ph{X.XX} & \ph{X.XX} & \ph{X.XX} & \ph{X.XX} \\
L5           & \ph{X.XX} & \ph{X.XX} & \ph{X.XX} & \ph{X.XX} & \ph{X.XX} & \ph{X.XX} \\
L6$^\dagger$ & \ph{X.XX} & 0.00      & \ph{X.XX} & N/A       & \ph{X.XX} & N/A       \\
Sacrum       & \ph{X.XX} & \ph{X.XX} & \ph{X.XX} & \ph{X.XX} & \ph{X.XX} & \ph{X.XX} \\
Left hip     & \ph{X.XX} & \ph{X.XX} & \ph{X.XX} & \ph{X.XX} & \ph{X.XX} & \ph{X.XX} \\
Right hip    & \ph{X.XX} & \ph{X.XX} & \ph{X.XX} & \ph{X.XX} & \ph{X.XX} & \ph{X.XX} \\
\bottomrule
\end{tabular}\\[2pt]
\footnotesize{$^\dagger$L6 restricted to lumbarization test cases ($n{=}3$).
TotalSegmentator L6 DSC\,=\,0.00 by architectural design (no L6 class).
N/A, not applicable.}
\end{table}

\subsection{Lumbarization Detection --- Primary Outcome}

\ph{On the 3 lumbarization test cases, SpineSurg-CT achieved mean L6 DSC of
X.XX (95\% CI: X.XX--X.XX), HD95 of X.XX\,mm. TotalSegmentator achieved L6
DSC\,=\,0.00 on all 3 cases ($p{<}0.001$).
Figure~\ref{fig:l6_example} shows a representative segmentation.}

\subsection{LSTV Classification}

\ph{Table~\ref{tab:classification} summarises LSTV classification performance.
Overall accuracy was X.XX.}

\begin{table}[ht]
\centering
\caption{LSTV Classification Performance}
\label{tab:classification}
\begin{tabular}{lccccc}
\toprule
System & Input & Sensitivity & Specificity & F1 & Subtype output \\
\midrule
SpineSurg-CT (this work) & CT & \ph{X.XX} & \ph{X.XX} & \ph{X.XX}
  & Lumb / Sacral / Normal \\
Kwak et al.~\citep{kwak2025} & Radiograph & 0.851 & 0.619 & N/A
  & Binary \\
\bottomrule
\end{tabular}\\[2pt]
\footnotesize{Kwak et al.\ comparison is binary (LSTV/Normal) on plain
radiograph; SpineSurg-CT provides three-class CT output. Direct comparison
should be interpreted cautiously given different modality and granularity.}
\end{table}

% =============================================================================
\section{Discussion}
% =============================================================================

SpineSurg-CT is the first CT segmentation model with explicit L6 supervision
and multi-class LSTV detection. The central finding is that SpineSurg-CT
achieves positive L6 DSC on lumbarization cases where TotalSegmentator scores
0.00 by architectural necessity. This is a capability gap, not a performance
margin.

The clinical motivation is direct. LSTV affects 4--18\% of the general
population~\citep{apazidis2011,nardo2012} and causes systematic vertebral
miscounting, the proximate cause of wrong-level spine surgery. Ishida et al.\
documented surgical level revision in 15\% of LSTV surgery
cases~\citep{ishida2017}; Albano et al.\ identified LSTV as the primary
anatomical risk factor in a systematic review of wrong-level
surgery~\citep{albano2024}. SpineSurg-CT removes the barrier of missing L6
representation in all existing automated CT tools.

The comparison with VERIDAH~\citep{moller2026} merits precise framing.
VERIDAH uses a 3D DenseNet with four learned classification heads and a
constrained sequence optimizer --- it does learn morphological features from CT
and is substantially more than a simple counting system. The key architectural
distinction from SpineSurg-CT is that VERIDAH represents L6 \textit{implicitly}:
during training, L6 labels are replaced with two consecutive L5 labels, and
the sequence optimizer detects lumbarization by finding that two consecutive L5
predictions minimise the global path cost. Consequently, there is no explicit L6
class, no direct voxel-level L6 supervision, and no 3D spatial output --- VERIDAH
produces vertebra centroid labels and requires pre-computed localizations from
an upstream segmentation model.

SpineSurg-CT is complementary: L6 is an explicit segmentation class learned
directly from annotated voxels, producing a 3D mask of the supernumerary segment
usable for surgical planning in a single end-to-end inference pass from raw CT.
Schinz et al.\ demonstrated at the thoracolumbar junction that explicit
shape-based classification achieves perfect agreement with functional anatomy
(nerve morphology, 100\% vs.\ 92.6\% for counting-based methods), while
counting-based classification produced discrepant labels in 13.3\% of
cases~\citep{schinz2025}. Although Schinz et al.\ addressed a different junction,
the principle that explicit morphological feature supervision produces more
reliable transitional vertebra detection applies to the lumbosacral problem.
The practical difference between VERIDAH and SpineSurg-CT is therefore implicit
sequence-level L6 detection with centroid output versus explicit voxel-supervised
L6 segmentation with a 3D spatial mask and probabilistic LSTV classification.

The registration-based sacralization synthesis pipeline represents a novel
methodological contribution. By performing ANTs SyN deformable registration
and field application entirely within the shared preprocessed $256^3$ unit
coordinate space, we eliminate the coordinate frame mismatch that arises when
fields computed in heterogeneous native NIfTI physical spaces are applied to
standardised training tensors. The resulting synthetic cases exhibit
simultaneous image and label deformation: the collapsed L5--S1 disc space,
enlarged transverse processes, and elongated sacrum morphology that characterise
radiological sacralization are present in the CT intensities, not merely in the
annotation. This is a qualitative advance over label-deletion approaches, which
remove the L5 annotation while leaving the underlying CT anatomy unchanged.
Restricting field application to the 23 training cases with complete pelvic
ground truth (sacrum and hip labels) ensures that every synthetic sacralization
sample contributes anatomically complete lumbosacral supervision signal,
including a fused L5+sacrum structure with intact hip labels.

\subsection{Limitations}

\textbf{Dataset size and LSTV case count.} The training cohort of 784 cases
with 25 LSTV cases is small relative to general-purpose segmentation benchmarks.
L6 DSC estimation on 3 lumbarization test cases and LSTV classification on
4 test cases has high variance; confidence intervals should be interpreted
accordingly and results treated as proof-of-concept pending external validation.

\textbf{Synthetic sacralization ceiling.} The matched-only restriction on
sacralization synthesis --- necessary to produce anatomically complete
sacrum and hip labels --- limits augmentation to the 23 pelvic-masked training
cases, yielding approximately 23 synthetic sacralization samples per epoch.
This represents the hard ceiling imposed by the intersection of the CTSpine1K
COLONOG and CTPelvic1K Dataset~1 cohorts. The model learns sacralized sacrum
morphology primarily from these 23 cases plus the 7 real sacralization training
cases, with limited diversity relative to the lumbarization augmentation
pipeline which can operate on all 510 normal training cases.

\textbf{Deformed-normal-sacrum approximation.} Synthetic sacralization is
generated by warping a morphologically normal matched reference toward the
sacralization appearance template. The resulting sacrum label covers the fused
L5+sacrum territory but reflects the geometry of a deformed normal sacrum rather
than the true bony morphology of sacralization, in which the L5 body is
incorporated into the superior sacral ala with altered trabecular architecture,
cortical thickening, and transverse process hypertrophy. The model may therefore
underestimate sacrum boundary precision at the L5 fusion zone. Sacralization
sacrum DSC on test cases should be interpreted with this approximation in mind.

\textbf{No pelvic ground truth for real sacralization training cases.} The 7
sacralization training cases have no matched CTPelvic1K pelvic mask, as the
CTSpine1K LSTV cases do not overlap with the CTPelvic1K COLONOG subset.
Consequently, sacrum and hip supervision for these cases derives entirely from
synthetic augmentation of pelvic-masked normals.

\textbf{Deformation field applied to sacrum label in synthetic sacralization.}
A more principled sacralization synthesis would warp only the lumbar vertebra
labels (L1--L5) of the normal reference toward the sacralization template while
holding the sacrum label fixed as the stationary fusion target. This would
ensure that the sacrum label remains correctly anchored to the underlying CT
anatomy and that the dilation-based L5 fusion step always operates against a
geometrically consistent sacrum boundary. In the current implementation,
\texttt{ants.apply\_transforms} is applied to the full label volume including
the sacrum class, which introduces a small but non-zero distortion of the sacrum
label relative to its true position. Although the 8-mm binary dilation used in
the fusion step provides some tolerance for this drift, the distortion may
reduce sacrum boundary precision at the L5 fusion zone.

Implementing the asymmetric approach requires knowledge of which voxels
constitute the sacrum in both the sacralization template images and the normal
reference volumes. For normal reference volumes this information is available
via the CTPelvic1K matched masks --- which is why synthesis is restricted to
the 23 pelvic-masked training cases. For the 7 real sacralization training cases,
however, CTSpine1K provides neither an L5 label (anatomically fused and therefore
unlabelled) nor a sacrum mask (no matched pelvic case exists). This creates a
fundamental labelling gap: it is not possible to separate the sacralized L5 from
the native sacrum in these cases without either (a)~manual expert annotation of
the L5--S1 fusion boundary, which is radiologically challenging and subjective
even for experienced spine radiologists, or (b)~pseudo-labelling using
TotalSegmentator~\citep{wasserthal2023}, which segments sacrum and hips by
morphological priors independently of vertebral counting and would correctly
identify the pelvic structures on all 784 cases. Pseudo-label quality on the
sacralization cases should be verified against the 33 manually annotated matched
cases before adoption. Either approach would enable the sacrum-fixed deformation
strategy described above and represents a concrete improvement target for a
subsequent study (see Section~\ref{sec:future}).

\textbf{No Castellvi subtype classification.} SpineSurg-CT detects LSTV type
but does not provide Castellvi subtype (Types~I--IV). This requires morphometric
analysis of transverse process height and sacral contact distance, planned as a
second-stage model.

\textbf{No explicit fused L5--S1 segmentation.} In sacralization, the fused L5
is incorporated into the sacrum label following CTSpine1K annotation convention.
The model does not distinguish the sacralized L5 from the native sacrum within
class 7.

\textbf{CT only; single-centre training data.} All training data originates
from COLONOG colonoscopy screening CTs acquired in a single imaging protocol.
External validation on diverse clinical CT acquisitions --- including different
scanners, reconstruction kernels, and patient populations --- is required before
deployment. MRI LSTV detection remains a separate problem addressed for the
labelling task by VERIDAH~\citep{moller2026}.

\subsection{Future Directions}
\label{sec:future}

\textbf{CycleGAN domain adaptation for sacralization synthesis.}
The primary limitation of the current registration-based approach is that the
synthetic image domain (geometrically deformed normal CT) differs from the real
sacralization image domain (true bony fusion with altered trabecular and cortical
architecture). Unpaired image-to-image translation using a CycleGAN
framework~\citep{zhu2017cyclegan} offers a principled approach to bridge this
domain gap: a generator trained on the 7 real sacralization cases would learn to
translate synthetically deformed CT intensities toward the appearance distribution
of real sacralization, without requiring paired data. This would substantially
improve the realism of synthetic sacralization cases beyond what geometric
registration alone can achieve.

\textbf{TotalSegmentator pseudo-labeling for expanded pelvic supervision.}
The matched-only synthesis ceiling (23 cases/epoch) is a direct consequence of
pelvic ground truth availability. TotalSegmentator~\citep{wasserthal2023}
segments sacrum and hips using morphological priors rather than vertebra counting
and will correctly identify the sacrum and hip structures on all 784 COLONOG
cases, including sacralization cases where the superior sacrum incorporates the
fused L5. Running TotalSegmentator pseudo-labeling on the full 784-case cohort
would provide sacrum and hip supervision for all normal training cases (510 cases
instead of 23), expanding the sacralization synthesis pool by $22\times$ and
eliminating the matched-only restriction. Pseudo-label quality should be verified
against the 33 manually annotated matched cases before adoption.

\textbf{Castellvi subtype classification.}
A second-stage model operating on the SpineSurg-CT segmentation output could
extract transverse process height and sacral contact distance to classify
Castellvi Types~I--IV. The spatial segmentation mask produced by SpineSurg-CT
provides the anatomical localization necessary for this morphometric analysis.

\textbf{External validation and prospective study.}
Validation on independent clinical CT cohorts, including pre-operative spine
surgery CTs with surgeon-confirmed vertebral level counts, is the critical next
step toward clinical deployment.

% =============================================================================
\section{Conclusion}
% =============================================================================

SpineSurg-CT is the first CT segmentation model with LSTV awareness built
into its architecture. It is the first to provide: (1)~a voxel-level 3D
segmentation mask for the supernumerary L6 vertebra under direct supervised
training; (2)~calibrated probabilistic LSTV subtype classification
(Normal/Lumbarization/Sacralization) from learned morphological features embedded
in the segmentation bottleneck; and (3)~both outputs simultaneously in a single
end-to-end inference pass from raw CT, without upstream vertebra localization or
counting steps. Existing tools partition these capabilities without unifying them:
TotalSegmentator segments without LSTV awareness; VERIDAH labels without
segmenting and requires an upstream segmentation model; Kwak et al.\ classify
without spatial output and require a separate imaging modality. The registration-based
sacralization synthesis pipeline --- performing ANTs SyN field computation and
application entirely in a shared preprocessed coordinate space and restricting
synthesis to pelvic-masked reference normals --- produces the first reported
simultaneous image-and-label sacralization augmentation for CT segmentation
training. Together these contributions represent a clinically meaningful advance
for preventing wrong-level spine surgery in patients with this anatomically
challenging and surgically consequential condition.

% =============================================================================
\section*{Acknowledgments}
% =============================================================================

The author thanks the Wayne State University HPC team for access to the warrior
cluster (H200 NVL GPU nodes, gmsap partition, allocation go2432). CTSpine1K and
CTPelvic1K datasets are publicly available and acknowledged per their respective
licenses.

% =============================================================================
\section*{Data and Code Availability}
% =============================================================================

CTSpine1K and CTPelvic1K are publicly available. Model weights, inference code,
train/val/test splits, and the CLAIM 2024 checklist will be released at
\url{[GitHub URL --- insert before submission]} under the Apache 2.0 license
upon publication.

% =============================================================================
\section*{Figure Captions}
% =============================================================================

\noindent\textbf{Figure 1.} Representative SpineSurg-CT segmentation output on
a lumbarization test case. Top row: axial slices at three levels showing the
predicted 10-class label overlay (L1--blue, L2--cyan, L3--teal, L4--green,
L5--lime, \textbf{L6--red}, sacrum--orange). Bottom row: corresponding ground
truth. The supernumerary L6 vertebra (red) is correctly identified between L5
and the sacrum; TotalSegmentator produces no L6 prediction on this case.
\ph{[Insert figure from inference visualization]}
\label{fig:l6_example}

\vspace{0.5em}
\noindent\textbf{Figure 2.} Registration-based sacralization synthesis pipeline.
Left column: original matched normal spine (CT image and label overlay). Centre
column: synthetic sacralization after ANTs SyN field application in preprocessed
$256^3$ unit space; note the collapsed L5--S1 disc space, enlarged transverse
processes, and fused L5+sacrum label (orange). Right column: a real sacralization
training case for comparison. Registration and field application are performed
entirely in the shared PIR $256^3$ 1\,mm isotropic coordinate space, eliminating
the coordinate frame mismatch inherent in native NIfTI-space registration pipelines.
\ph{[Insert from augmentation QC visuals]}

% ─────────────────────────────────────────────────────────────────────────────
\bibliographystyle{unsrt}
\bibliography{SpineSurgCT_references}

\end{document}
